\begin{statementbox}
	\textbf{\textcolor{CStatement}{条件}}
	\begin{description}[style=nextline, leftmargin=2.8em, labelsep=0.5em, itemsep=0.35em]
		\item[\textbf{\textcolor{CStatement}{条件 1（角度任意）：}}] $\alpha$ 为任意实数，即 $\alpha\in\mathbb{R}$。
		\item[\textbf{\textcolor{CStatement}{条件 2（余弦非零）：}}] 对商数关系 $\tan\alpha=\dfrac{\sin\alpha}{\cos\alpha}$，要求 $\cos\alpha\neq0$，即 $\alpha\neq\dfrac{\pi}{2}+k\pi\ (k\in\mathbb{Z})$。
	\end{description}

	\textbf{\textcolor{CStatement}{结论}}
	\begin{description}[style=nextline, leftmargin=2.8em, labelsep=0.5em, itemsep=0.45em]
		\item[\textbf{\textcolor{CStatement}{结论一（平方关系）：}}]
			\textbf{\textcolor{CStatement}{直观描述：}}
			任意角的正弦平方与余弦平方之和恒为1。
			\[
				\sin^2\alpha + \cos^2\alpha = 1.
			\]

		\item[\textbf{\textcolor{CStatement}{结论二（商数关系）：}}]
			\textbf{\textcolor{CStatement}{直观描述：}}
			当余弦不为零时，正切等于正弦除以余弦。
			\[
				\tan\alpha = \frac{\sin\alpha}{\cos\alpha}\quad (\cos\alpha \neq 0).
			\]

		\item[\textbf{\textcolor{CStatement}{推广结论（平方变形）：}}]
			\textbf{\textcolor{CStatement}{直观描述：}}
			由平方关系直接移项或与和差平方公式联立。
			\[
				\sin^2\alpha = 1 - \cos^2\alpha,
				\qquad
				\cos^2\alpha = 1 - \sin^2\alpha.
			\]
			\[
				(\sin\alpha \pm \cos\alpha)^2 = 1 \pm 2\sin\alpha\cos\alpha.
			\]
	\end{description}
\end{statementbox}
